\documentclass[11pt]{article}
\usepackage[T1]{fontenc}
\usepackage{lmodern}
\usepackage[margin=1in]{geometry}
\usepackage{amsmath,amssymb,amsthm,booktabs,microtype,listings}
\usepackage[hidelinks]{hyperref}
\hypersetup{pdftitle={A zero-free disk for the minimum cost of a four-edge bipartite matching},
  pdfauthor={Xiaojian Zeng}}
\newtheorem{theorem}{Theorem}
\newtheorem{lemma}[theorem]{Lemma}
\newcommand{\E}{\mathbb E}
\newcommand{\R}{\mathbb R}
\newcommand{\C}{\mathbb C}
\DeclareMathOperator{\Ree}{Re}
\title{A zero-free disk for the minimum cost\\of a four-edge bipartite matching}
\author{Xiaojian Zeng\\[3pt]
  \small PBC School of Finance, Tsinghua University, 100190\\
  \small\texttt{xiaojianzeng@foxmail.com}}
\date{September 2026}
\begin{document}
\maketitle
\begin{abstract}
Let $C_{4,m,n}$ be the minimum total cost of a four-edge matching in
$K_{m,n}$, with independent exponential edge costs of mean one.
For every pair of integers $m,n\ge4$, we prove that
$\E e^{tC_{4,m,n}}$ has no zero in $|t|<mn/4$.
Its reduced rational numerator has no zero even in the closed disk.
The proof establishes a stronger algebraic statement for all real
$m,n\ge4$: a M\"obius transform of the explicit sextic numerator has
positive Hurwitz determinants, certified by exact integer coefficients
after shifting $m$ and $n$ by four. The unrestricted matching-size
conjecture of W\"astlund remains unresolved.
\end{abstract}

\section{The four-edge statement}

Assign independent rate-one exponential variables $X_{ij}$ to the edges
of the complete bipartite graph $K_{m,n}$. A matching is a set of edges
with no shared endpoints. For $1\le k\le\min(m,n)$, put
\[
 C_{k,m,n}=\min_{M:\,|M|=k}\sum_{(i,j)\in M}X_{ij},
 \qquad F_{k,m,n}(t)=\E e^{tC_{k,m,n}}.
\]
W\"astlund conjectures that $F_{k,m,n}$ is nonzero for
$|t|<mn/k$ \cite{wastlund2026} (Conjecture 11). We settle the complete
four-edge range of that conjecture.

\begin{theorem}\label{thm:main}
For every pair of integers $m,n\ge4$,
\[
 F_{4,m,n}(t)\neq0\qquad(t\in\C,\ |t|<mn/4).
\]
Every zero of the reduced numerator of its rational continuation has
modulus strictly greater than $mn/4$. More strongly, the sextic
$P_{m,n}$ defined in \eqref{eq:coefficients} has no root in
$|t|\le mn/4$ for every pair of real numbers $m,n\ge4$.
\end{theorem}

The distinction between the open and closed disks is necessary:
$t=mn/4$ is a pole of the rational continuation. The theorem makes
no assertion of a finite value there. Nonintegral $m,n$ occur only
in the algebraic part of the statement.

The closest zero-exclusion results in \cite{wastlund2026} (Section 4.3 and
Propositions 15--16) cover all $k\le3$ and leave at most
finitely many possible exceptional pairs $(m,n)$ for each fixed $k$;
the latter result gives no explicit cutoff. Theorem~\ref{thm:main}
removes all possible exceptions for $k=4$.
Alm and Sorkin had already developed exact distribution calculations,
including Laplace transforms, for $k\le4$ \cite{almsorkin2002}.
Thus our contribution is the uniform zero exclusion. The rational
recursion and the ability to compute a four-edge transform are prior work.
No equivalent four-edge zero-exclusion theorem was identified in the
primary sources inspected through 14 September 2026; this is a bounded
literature comparison, not a claim of exhaustive priority.

\section{The rational numerator}

Write $r=mn$ and $s=m+n$, and set
\begin{equation}\label{eq:AQ}
 A=\prod_{j=0}^{3}(m-j)(n-j),\qquad
 Q(t)=\prod_{\substack{i,j\ge0\\ i+j<4}}
       \bigl((m-i)(n-j)-(4-i-j)t\bigr).
\end{equation}
Define $P(t)=P_{m,n}(t)=\sum_{a=0}^6p_a t^a$ by
\begin{equation}\label{eq:coefficients}
\begin{aligned}
 p_0={}&r^3(r-s+1)^2(r-2s+4),\\
 p_1={}&-r^2(r-s+1)(12s^2-25sr-28s+10r^2+37r+12),\\
 p_2={}&r\bigl(40r^3-117sr^2+97s^2r-22s^3+129r^2\\
       &\qquad\qquad{}-179sr+44s^2+75r-16s-8\bigr),\\
 p_3={}&12s^3-107s^2r+187sr^2+117sr+4s\\
       &\qquad\qquad{}-82r^3-145r^2-58r-8,\\
 p_4={}&42s^2-143sr-2s+91r^2+47r+8,\\
 p_5={}&42s-52r+10,\qquad p_6=12.
\end{aligned}
\end{equation}

\begin{lemma}\label{lem:rational}
For integers $m,n\ge4$, the rational continuation of the MGF is
\begin{equation}\label{eq:rational}
 F_{4,m,n}(t)=\frac{A P(t)}{Q(t)}.
\end{equation}
\end{lemma}

\begin{proof}
We give the finite calculation from W\"astlund's recursion
\cite{wastlund2026} (Theorem 8 and equation (9)), including its separate
rule for the ordinary MGF. Let $N_i=n-4+i$ and let
$L_{i,j}=F^{(i-j)}_{i,m,N_i}$ in that source's ghost-function notation.
Only the diagonal functions are ordinary MGFs.
Start with $L_{0,0}=1$. For $1\le i\le4$ and $0\le j<i$, use
\begin{equation}\label{eq:recurrence}
 L_{i,j}=\frac{N_iL_{i-1,j}+(m-j)(i-j+1)tL_{i,j-1}
                         +(m-j)N_iL_{i-1,j-1}}
                   {(m-j)N_i-(i-j)t},
\end{equation}
where out-of-range terms are zero. On the diagonal use
\begin{equation}\label{eq:diagonal}
 L_{i,i}=\frac{t}{N_i}L_{i,i-1}+L_{i-1,i-1}.
\end{equation}
Increasing $i$, and then $j$, determines all fifteen cells. Combining
these fractions gives $L_{4,4}=AP/Q$, with the coefficients
\eqref{eq:coefficients}. This is a rational-function identity in
$\mathbb Q(m,n,t)$, obtained by clearing denominators; it also gives
$AP(0)=Q(0)$ and $\deg_tP=6$.
The accompanying program \texttt{reconstruct.py} checks that identity
directly from \eqref{eq:recurrence}--\eqref{eq:diagonal} and the displayed
coefficients.

For $m,n\ge4$, $N_i\ge1$ and $m-j\ge1$ whenever $j<i$.
The diagonal rule never divides by $m-i$, so the endpoint $m=i=4$
is included. Each recurrence denominator is a nonzero polynomial in
$t$ after specialization. Consequently the generic rational identity
specializes throughout the asserted domain.
\end{proof}

\section{An exact Hurwitz certificate}

Let $R=mn/4$ and define
\begin{equation}\label{eq:H}
 H(z)=\frac{4^6}{(mn)^3}(1+z)^6
       P\!\left(R\frac{1-z}{1+z}\right)
      =h_0z^6+h_1z^5+\cdots+h_6.
\end{equation}
The left side is interpreted as a polynomial after clearing the
substitution denominator. Set $x=m-4$, $y=n-4$, and
$v(u)=(1,u,u^2,u^3)^{\mathsf T}$. Direct expansion of
\eqref{eq:coefficients} gives
\begin{equation}\label{eq:matrices}
 h_a=v(x)^{\mathsf T}B_a v(y),\qquad 0\le a\le6,
\end{equation}
where all seven integer matrices $B_a$ are printed in
Appendix~\ref{app:matrices}.

Form the Hurwitz matrix
\begin{equation}\label{eq:hurwitz}
 \mathcal H=\left(\begin{array}{cccccc}
 h_1&h_3&h_5&0&0&0\\
 h_0&h_2&h_4&h_6&0&0\\
 0&h_1&h_3&h_5&0&0\\
 0&h_0&h_2&h_4&h_6&0\\
 0&0&h_1&h_3&h_5&0\\
 0&0&h_0&h_2&h_4&h_6
 \end{array}\right),
\end{equation}
and denote its leading $j$ by $j$ determinant by $\Delta_j$.

\begin{lemma}\label{lem:positive}
Every $h_a$ has strictly positive coefficients on the full support
$\{x^uy^v:0\le u,v\le3\}$. For $1\le j\le6$, every coefficient
of $\Delta_j$ on the full support
$\{x^uy^v:0\le u,v\le3j\}$ is strictly positive, and there are
no other terms.
\end{lemma}

\begin{proof}
The first assertion is visible in the matrices $B_a$.
Expand the determinants of \eqref{eq:hurwitz} over $\mathbb Z[x,y]$.
The exact coefficient counts and minima are as follows.
\begin{center}
\begin{tabular}{crrr}
$j$ & Degree in each variable & Number of terms & Least coefficient\\
1&3 &16 &82200\\
2&6 &49 &5637588000\\
3&9 &100&176698647552000\\
4&12&169&1539481522876416000\\
5&15&256&2195347272420556800000\\
6&18&361&237097505421420134400000\\
\end{tabular}
\end{center}
Appendix~\ref{app:arithmetic} specifies the entire expansion using
integer additions and multiplications, and checks every support entry
and coefficient. Thus the table summarizes a finite algebraic
certificate, with no parameter sampling or root approximation.
Expansion along the last column also gives
$\Delta_6=h_6\Delta_5$.
\end{proof}

For $x,y\ge0$, Lemma~\ref{lem:positive} implies $h_0>0$ and
$\Delta_j>0$ for all six $j$. The positive constant terms ensure
strictness on both axes, including $x=y=0$.
The real Routh--Hurwitz criterion therefore puts every root of $H$
in $\Ree z<0$; see Holtz \cite{holtz2003} (Theorem 3 and its proof).
To reconcile coefficient conventions explicitly, apply that source's
ascending-coefficient form to
$f(w)=w^6H(1/w)=h_0+h_1w+\cdots+h_6w^6$.
Its first seven leading minors are
$h_0,h_0\Delta_1,\ldots,h_0\Delta_6$.
All are positive. Since $h_0,h_6>0$, inversion bijects the roots of
$f$ and $H$, and $\Ree(1/w)=\Ree(w)/|w|^2$ preserves their half-plane.

Exact coefficient positivity is useful for certifying a whole parameter
domain from finite data. A related implementation appears in
Zeng \cite{zeng2026lens} (Theorem 10, equations (8.25)--(8.27)),
where rational Bernstein coefficients certify a two-parameter
spectral-set inequality on a box. Here shifted monomial coefficients
certify Hurwitz determinants on an unbounded quadrant. The
spectral-set result supplies a methodological comparison, not a
matching-cost theorem.

\section{Proof of the zero-free disk}

\begin{proof}[Proof of Theorem~\ref{thm:main}]
First let $m,n\ge4$ be real. For $|t|\le R$ and $t\neq-R$, put
$z=(R-t)/(R+t)$. Then
\[
 \Ree z=\frac{R^2-|t|^2}{|R+t|^2}\ge0,
 \qquad 1+z=\frac{2R}{R+t}\neq0.
\]
Strict Hurwitz stability and \eqref{eq:H} imply $P(t)\neq0$.
The exceptional point is covered by the leading coefficient:
\[
 h_0=\frac{4^6}{(mn)^3}P(-R)>0.
\]
This proves the algebraic assertion on the closed disk.

All roots of $Q$ are positive real numbers
$\rho_{ij}=(m-i)(n-j)/(4-i-j)$, where $i+j<4$. Since
\[
 \frac{(m-i)(n-j)}{mn}
 =\left(1-\frac{i}{m}\right)\left(1-\frac{j}{n}\right)
 \ge\left(1-\frac{i}{4}\right)\left(1-\frac{j}{4}\right)
 \ge1-\frac{i+j}{4},
\]
we have $\rho_{ij}\ge R$. Also $A>0$.
For integer parameters, \eqref{eq:rational} is therefore nonzero
throughout the open disk. After any such specialization,
cancellation of common polynomial factors can only remove roots
from $P$, which has leading coefficient $12$ and is never the zero
polynomial. This proves the reduced-numerator assertion.
The factor indexed by $(i,j)=(0,0)$ vanishes at $t=R$, whereas
$P(R)=(mn)^3h_6/4^6>0$. Thus $R$ is indeed a pole.

Finally, the rational expression represents the analytic expectation
on the claimed disk. For every fixed nonnegative cost array, the
minimum cost is at most the average over all four-edge matchings.
Each edge occurs in a fraction $4/(mn)$ of those matchings, so
\[
 C_{4,m,n}\le\frac{4}{mn}\sum_{i,j}X_{ij},\qquad
 \E e^{aC_{4,m,n}}\le
       \left(1-\frac{4a}{mn}\right)^{-mn}<\infty
       \quad(0\le a<R).
\]
For each compact subset of $|t|<R$, choose $a<R$ strictly larger
than the real parts on that set. Domination by $e^{aC_{4,m,n}}$
gives analyticity. The source identity agrees with the expectation
near zero, hence throughout the disk by analytic continuation.
\end{proof}

The result concerns exactly four edges. It proves no case with
$k\ge5$ and gives no Gaussian scaling limit when the matching size
grows. The useful finite step is the positivity of all six Hurwitz
determinants on the entire shifted quadrant; positivity of the
seven coefficients of $H$ alone would not suffice.

\paragraph{Computation and assistance.}
The rational identity and the transformed coefficients were checked
with SymPy 1.14.0 \cite{meurer2017}; a separate implementation using
only Python integer arithmetic rebuilt the cleared-denominator
recursion and all determinant coefficients. Both the displayed
matrices and the complete support rectangles were checked.
The accompanying \texttt{reconstruct.py} checks the manuscript's
sextic and its transform, and \texttt{hurwitz\_certificate.py}
reproduces the determinant table without third-party packages.
These are exact symbolic computations over an infinite parameter
domain, not a finite sweep. OpenAI Codex assisted with proof
development, programming, source checking, and manuscript drafting.
A separate AI review pass recomputed the rational identity and
certificate; this is neither human peer review nor proof-assistant
formalization. The mathematical argument consists of the stated
identities, exact sign certificates, and the classical stability
criterion.

\appendix
\section{The transformed coefficients}\label{app:matrices}

In \eqref{eq:matrices}, rows and columns are indexed by $0,1,2,3$;
the $(u,v)$ entry of $B_a$ is $[x^uy^v]h_a$.
The following matrices contain all $112$ coefficients.
\small
\[
B_0=\left(\begin{array}{cccc}
20519424 & 21197056 & 7200000 & 806400 \\
21197056 & 21858816 & 7416032 & 829920 \\
7200000 & 7416032 & 2514760 & 281400 \\
806400 & 829920 & 281400 & 31500
\end{array}\right)
\]
\[
B_1=\left(\begin{array}{cccc}
33546240 & 38182912 & 14102784 & 1701120 \\
38182912 & 43108608 & 15818560 & 1897792 \\
14102784 & 15818560 & 5775296 & 690080 \\
1701120 & 1897792 & 690080 & 82200
\end{array}\right)
\]
\[
B_2=\left(\begin{array}{cccc}
22417920 & 27708672 & 11058688 & 1434368 \\
27708672 & 34036992 & 13489952 & 1737504 \\
11058688 & 13489952 & 5309384 & 679384 \\
1434368 & 1737504 & 679384 & 86420
\end{array}\right)
\]
\[
B_3=\left(\begin{array}{cccc}
7274496 & 10035200 & 4387328 & 616960 \\
10035200 & 13540864 & 5832576 & 810880 \\
4387328 & 5832576 & 2482112 & 341504 \\
616960 & 810880 & 341504 & 46544
\end{array}\right)
\]
\[
B_4=\left(\begin{array}{cccc}
1101312 & 1815296 & 907520 & 142336 \\
1815296 & 2783232 & 1329440 & 202272 \\
907520 & 1329440 & 616280 & 91752 \\
142336 & 202272 & 91752 & 13428
\end{array}\right)
\]
\[
B_5=\left(\begin{array}{cccc}
73728 & 147456 & 89344 & 16640 \\
147456 & 268544 & 147776 & 25408 \\
89344 & 147776 & 75904 & 12384 \\
16640 & 25408 & 12384 & 1944
\end{array}\right)
\]
\[
B_6=\left(\begin{array}{cccc}
1536 & 3840 & 3072 & 768 \\
3840 & 8448 & 5856 & 1248 \\
3072 & 5856 & 3480 & 648 \\
768 & 1248 & 648 & 108
\end{array}\right)
\]
\normalsize

\section{Reproducing the determinant certificate}\label{app:arithmetic}

The code below uses a dictionary whose entry \texttt{(u,v):c}
represents $cx^uy^v$. Set \texttt{B} to the list
$[B_0,\ldots,B_6]$ in Appendix~\ref{app:matrices}.
The accompanying file \texttt{hurwitz\_certificate.py} includes those
matrices and runs as \texttt{python3 hurwitz\_certificate.py}.
The determinant recursion is ordinary first-row Laplace expansion.
It compares the entire support rectangle, the coefficient minimum,
and the identity $\Delta_6=h_6\Delta_5$.
\begin{lstlisting}[language=Python,basicstyle=\ttfamily\footnotesize,
  columns=fullflexible,keepspaces=true,showstringspaces=false,
  breaklines=true,frame=single,framesep=5pt]
def add(p, q, sign=1):
    out = p.copy()
    for e, c in q.items():
        out[e] = out.get(e, 0) + sign*c
    return {e: c for e, c in out.items() if c}

def mul(p, q):
    out = {}
    for (a, b), c in p.items():
        for (u, v), d in q.items():
            e = (a+u, b+v)
            out[e] = out.get(e, 0) + c*d
    return {e: c for e, c in out.items() if c}

def det(M):
    if not M:
        return {(0, 0): 1}
    out = {}
    for j, entry in enumerate(M[0]):
        if entry:
            sub = [row[:j] + row[j+1:] for row in M[1:]]
            out = add(out, mul(entry, det(sub)), (-1)**j)
    return out

def certificate():
    h = [{(u, v): b[u][v] for u in range(4)
          for v in range(4)} for b in B]
    assert all(c > 0 for p in h for c in p.values())
    H = [[h[2*j-i+1] if 0 <= 2*j-i+1 <= 6 else {}
          for j in range(6)] for i in range(6)]
    minima = [82200, 5637588000, 176698647552000,
              1539481522876416000, 2195347272420556800000,
              237097505421420134400000]
    D = []
    for j in range(1, 7):
        p = det([row[:j] for row in H[:j]])
        support = {(u, v) for u in range(3*j+1)
                   for v in range(3*j+1)}
        assert set(p) == support
        assert min(p.values()) == minima[j-1] > 0
        D.append(p)
        print(j, len(p), min(p.values()))
    assert D[5] == mul(h[6], D[4])
    return h, D

if __name__ == "__main__":
    certificate()
\end{lstlisting}

\clearpage
\bibliographystyle{plain}
\bibliography{references}
\end{document}
